% ========================================================
% BÀI TẬP TOÁN 5 - KẾT NỐI TRI THỨC VỚI CUỘC SỐNG
% Bài 15: Ki-lô-mét vuông. Héc-ta
% Định dạng: BTTeX Part III - Trắc nghiệm điền số ngắn
% ========================================================

% Câu 1
\begin{ex}
  Một ki-lô-mét vuông bằng bao nhiêu héc-ta?
  \shortans{100}
  \loigiai{
    $1\text{ km}^2 = 100\text{ ha}$.
  }
\end{ex}

% Câu 2
\begin{ex}
  Một héc-ta bằng bao nhiêu mét vuông?
  \shortans{10000}
  \loigiai{
    $1\text{ ha} = 10\,000\text{ m}^2$.
  }
\end{ex}

% Câu 3
\begin{ex}
  Đổi $5\text{ km}^2$ thành bao nhiêu héc-ta?
  \shortans{500}
  \loigiai{
    $5\text{ km}^2 = 5 \times 100 = 500\text{ ha}$.
  }
\end{ex}

% Câu 4
\begin{ex}
  Đổi $12\text{ ha}$ thành bao nhiêu mét vuông?
  \shortans{120000}
  \loigiai{
    $12\text{ ha} = 12 \times 10\,000 = 120\,000\text{ m}^2$.
  }
\end{ex}

% Câu 5
\begin{ex}
  Đổi $400\text{ ha}$ thành bao nhiêu ki-lô-mét vuông?
  \shortans{4}
  \loigiai{
    $400\text{ ha} = 400 : 100 = 4\text{ km}^2$.
  }
\end{ex}

% Câu 6
\begin{ex}
  Đổi $80\,000\text{ m}^2$ thành bao nhiêu héc-ta?
  \shortans{8}
  \loigiai{
    $80\,000\text{ m}^2 = 80\,000 : 10\,000 = 8\text{ ha}$.
  }
\end{ex}

% Câu 7
\begin{ex}
  Đổi $3\text{ km}^2\text{ } 25\text{ ha}$ thành bao nhiêu héc-ta?
  \shortans{325}
  \loigiai{
    $3\text{ km}^2 = 300\text{ ha}$. Do đó $3\text{ km}^2\text{ } 25\text{ ha} = 325\text{ ha}$.
  }
\end{ex}

% Câu 8
\begin{ex}
  Đổi $2\text{ ha } 500\text{ m}^2$ thành bao nhiêu mét vuông?
  \shortans{20500}
  \loigiai{
    $2\text{ ha} = 20\,000\text{ m}^2$. Do đó $2\text{ ha } 500\text{ m}^2 = 20\,500\text{ m}^2$.
  }
\end{ex}

% Câu 9
\begin{ex}
  Một khu rừng hình chữ nhật có chiều dài $3\text{ km}$, chiều rộng $2\text{ km}$. Diện tích khu rừng đó là bao nhiêu héc-ta?
  \shortans{600}
  \loigiai{
    Diện tích là $3 \times 2 = 6\text{ km}^2 = 600\text{ ha}$.
  }
\end{ex}

% Câu 10
\begin{ex}
  Một nông trường hình chữ nhật có chiều dài $500\text{ m}$, chiều rộng $400\text{ m}$. Diện tích nông trường đó là bao nhiêu héc-ta?
  \shortans{20}
  \loigiai{
    Diện tích là $500 \times 400 = 200\,000\text{ m}^2 = 200\,000 : 10\,000 = 20\text{ ha}$.
  }
\end{ex}

% Câu 11
\begin{ex}
  Một hồ nước hình vuông có cạnh dài $1\,000\text{ m}$. Diện tích hồ nước đó là bao nhiêu ki-lô-mét vuông?
  \shortans{1}
  \loigiai{
    $1\,000\text{ m} = 1\text{ km}$. Diện tích là $1 \times 1 = 1\text{ km}^2$.
  }
\end{ex}

% Câu 12
\begin{ex}
  Một cánh đồng lúa có diện tích $5\text{ ha}$. Người ta đã gặt được $\frac{2}{5}$ diện tích. Hỏi diện tích lúa còn lại chưa gặt là bao nhiêu héc-ta?
  \shortans{3}
  \loigiai{
    Diện tích còn lại là: $5 \times (1 - \frac{2}{5}) = 5 \times \frac{3}{5} = 3\text{ ha}$.
  }
\end{ex}

% Câu 13
\begin{ex}
  Một khu công nghiệp có diện tích $150\text{ ha}$. Trong đó diện tích xây dựng nhà xưởng chiếm $\frac{2}{3}$ diện tích. Diện tích xây dựng nhà xưởng là bao nhiêu héc-ta?
  \shortans{100}
  \loigiai{
    Diện tích nhà xưởng là: $150 \times \frac{2}{3} = 100\text{ ha}$.
  }
\end{ex}

% Câu 14
\begin{ex}
  Người ta cấy lúa trên một thửa ruộng rộng $2\text{ ha}$. Trung bình mỗi héc-ta thu hoạch được $60\text{ tạ}$ thóc. Hỏi cả thửa ruộng thu hoạch được bao nhiêu tấn thóc?
  \shortans{12}
  \loigiai{
    Tổng số thóc là: $60 \times 2 = 120\text{ tạ} = 12\text{ tấn}$.
  }
\end{ex}

% Câu 15
\begin{ex}
  Một khu sinh thái rộng $2\text{ km}^2$. Người ta dành $\frac{1}{4}$ diện tích để làm hồ nước. Diện tích hồ nước là bao nhiêu héc-ta?
  \shortans{50}
  \loigiai{
    $2\text{ km}^2 = 200\text{ ha}$. Diện tích hồ nước là $200 \times \frac{1}{4} = 50\text{ ha}$.
  }
\end{ex}

% Câu 16
\begin{ex}
  Diện tích nước Việt Nam khoảng $331\,212\text{ km}^2$. Chữ số hàng trăm nghìn của số này là bao nhiêu?
  \shortans{3}
  \loigiai{
    Chữ số ở hàng trăm nghìn là 3.
  }
\end{ex}

% Câu 17
\begin{ex}
  Đổi $0{,}5\text{ km}^2$ thành bao nhiêu héc-ta?
  \shortans{50}
  \loigiai{
    $0{,}5\text{ km}^2 = 0{,}5 \times 100 = 50\text{ ha}$.
  }
\end{ex}

% Câu 18
\begin{ex}
  Đổi $0{,}25\text{ ha}$ thành bao nhiêu mét vuông?
  \shortans{2500}
  \loigiai{
    $0{,}25\text{ ha} = 0{,}25 \times 10\,000 = 2\,500\text{ m}^2$.
  }
\end{ex}

% Câu 19
\begin{ex}
  Một mảnh đất hình chữ nhật có chu vi $1\,400\text{ m}$, chiều rộng bằng $\frac{3}{4}$ chiều dài. Diện tích mảnh đất là bao nhiêu héc-ta?
  \shortans{12}
  \loigiai{
    Nửa chu vi là $700\text{ m}$. Chiều dài là $700 : (3 + 4) \times 4 = 400\text{ m}$. Chiều rộng là $300\text{ m}$. Diện tích là $400 \times 300 = 120\,000\text{ m}^2 = 12\text{ ha}$.
  }
\end{ex}

% Câu 20
\begin{ex}
  Một trang trại nuôi bò có diện tích cỏ là $40\text{ ha}$. Mỗi héc-ta cỏ đủ nuôi $5$ con bò. Trang trại đó nuôi được bao nhiêu con bò?
  \shortans{200}
  \loigiai{
    Số con bò nuôi được là: $40 \times 5 = 200$ (con).
  }
\end{ex}

% Câu 21
\begin{ex}
  Đổi $1\,500\,000\text{ m}^2$ thành bao nhiêu héc-ta?
  \shortans{150}
  \loigiai{
    $1\,500\,000 : 10\,000 = 150\text{ ha}$.
  }
\end{ex}

% Câu 22
\begin{ex}
  Đổi $7\text{ km}^2$ thành bao nhiêu mét vuông? (Điền số nghìn mét vuông).
  \shortans{7000}
  \loigiai{
    $7\text{ km}^2 = 7\,000\,000\text{ m}^2 = 7\,000$ nghìn mét vuông.
  }
\end{ex}

% Câu 23
\begin{ex}
  Một cánh rừng có diện tích $8\text{ km}^2$. Sau một vụ trồng cây diện tích tăng thêm $50\text{ ha}$. Tổng diện tích rừng lúc này là bao nhiêu héc-ta?
  \shortans{850}
  \loigiai{
    $8\text{ km}^2 = 800\text{ ha}$. Tổng diện tích là $800 + 50 = 850\text{ ha}$.
  }
\end{ex}

% Câu 24
\begin{ex}
  So sánh $2\text{ km}^2$ và $250\text{ ha}$. Số đo lớn hơn nhiều hơn số đo bé bao nhiêu héc-ta?
  \shortans{50}
  \loigiai{
    $2\text{ km}^2 = 200\text{ ha}$. Vì $250\text{ ha} > 200\text{ ha}$ nên hiệu là $250 - 200 = 50\text{ ha}$.
  }
\end{ex}

% Câu 25
\begin{ex}
  Một công viên hình chữ nhật dài $800\text{ m}$, rộng $500\text{ m}$. Diện tích công viên là bao nhiêu héc-ta?
  \shortans{40}
  \loigiai{
    Diện tích là $800 \times 500 = 400\,000\text{ m}^2 = 40\text{ ha}$.
  }
\end{ex}
